\chapter{Научные основы коннектомного моделирования и измерения ECAP}
\label{chap:foundations}

Коннектомное моделирование и расчёт вызванного электрического ответа решают разные задачи нейроинженерии. В первом случае исследуется, как сведения о связях между нейронами ограничивают возможную динамику сети. Во втором требуется установить, какой сигнал возникает на регистрирующих электродах при заданном воздействии на нервную ткань. Для их сопряжения необходимо определить переносимый объект и показать, какое свойство исходной модели имеет значение для целевой измерительной задачи. Само наличие подробного анатомического графа такого обоснования не даёт \cite{s743,s285}.

В главе последовательно рассматриваются наблюдаемые реакции \textit{Drosophila melanogaster}, коннектомные данные, механизмы формирования вызванного составного потенциала действия и методы проверки моделей. Литературное сопоставление ограничено вопросами, необходимыми для постановки исследования: совместимостью структуры и функциональных наблюдений, выбором представления динамики, достаточностью физического описания регистрации и независимостью сравнительной оценки. Результаты рассмотренных работ отделяются от требований к дальнейшему исследованию. Представленные ниже схемы обобщают эти требования; экспериментальные результаты предлагаемого переноса в главе не приводятся.

Обзор построен как критическое тематическое сопоставление, а не как систематическая оценка всей литературы по нейронному моделированию. В него включены работы, позволяющие разобрать отдельные звенья постановки: регистрацию ответа Drosophila, ограничения анатомического графа, получение динамики, электрическое возбуждение и регистрацию ECAP, а также обучение на синтетических данных. При выборе источников предпочтение отдаётся первичным исследованиям с описанием наблюдаемой величины, исходных данных и проверки результата. Теоретические работы и патентные описания используются в соответствии с их назначением; наличие формальной модели или опубликованного патента не рассматривается как свидетельство экспериментальной эффективности.

Работы сопоставляются по объекту моделирования, измеряемому выходу, условиям настройки и независимой проверки. Для препринтов существенна версия текста; для частично доступных публикаций вывод ограничивается доступными сведениями. Сходство архитектур или терминов рассматривается как основание для сравнения методов, но не заменяет совпадения исследовательских задач. Такой подход позволяет выделить обоснованные компоненты будущего метода и установить, какие предположения требуют самостоятельной проверки, не делая вывода о приоритете на основании ограниченного отбора литературы.

\section{Физический стимул, нейронная реакция и защитное поведение Drosophila}
\label{sec:constructs}

\subsection{Наблюдаемые величины и способы исследования ноцицепции}

Исследование ноцицепции начинается с определения воздействия и наблюдаемого ответа. Температура поверхности, механическая сила, длительность контакта и положение зоны стимуляции описывают физический вход. Мембранный ток, частота разрядов и изменение флуоресценции характеризуют различные стороны нейронной активности. Поворот, сокращение тела, перекатывание и избегание относятся к поведению целого организма. Между этими уровнями существуют причинные связи, однако эксперимент на одном из них не заменяет наблюдений на остальных \cite{s778,s810,s031}.

Работа о гене \textit{painless} связала изменения защитной реакции личинок с молекулярным механизмом чувствительности к повреждающим воздействиям. В ней исследовались как периферические нервные ответы, так и характерное перекатывание личинки. Эти показатели имеют разную размерность и разную единицу наблюдения: электрическая активность относится к регистрируемой нервной структуре, а двигательная реакция --- к животному в условиях конкретной пробы. Их совместное использование усиливает механистическую интерпретацию, но не позволяет считать поведение прямой записью нейронного сигнала \cite{s778}.

Для механической ноцицепции существенны результаты исследований Piezo и Pickpocket. В работе о Piezo электрофизиологические токи получали у диссоциированных личиночных нейронов при механическом воздействии через регистрирующую пипетку; поведенческие испытания составляли отдельную часть исследования. При переносе этих данных в модель следует сохранять различие между воздействием на изолированную клетку и воздействием на тело животного. Работа о Pickpocket прежде всего устанавливает связь молекулярной манипуляции с защитной реакцией. Она не предоставляет непрерывной записи вызванной активности, которую можно непосредственно использовать как целевой временной ряд \cite{s810,s795}.

Модальность стимула также определяет границы сопоставления. Исследования холодовой ноцицепции объединяют сведения о связности, измерения нейронной активности и поведенческие пробы. Оптогенетическая активация, подавление клеточных типов и регистрация кальциевых показателей проверяют разные звенья предполагаемого механизма. Холодовое сокращение тела и тепловое перекатывание нельзя объединять под одной выходной меткой без обоснования того, какое свойство реакции должно сохраняться \cite{s031,s901}.

Кальциевое изображение требует отдельного описания наблюдения. Регистрируемая флуоресценция зависит от внутриклеточных процессов, свойств индикатора, частоты съёмки и обработки. Её временной ход отличается от последовательности потенциалов действия. Если функциональные данные представлены флуоресценцией, сравнение с расчётной сетью должно выполняться после соответствующего преобразования состояния модели. Прямое приравнивание отсчёта флуоресценции к мгновенной частоте разрядов вводит дополнительное, часто не проверенное допущение \cite{s212,s285}.

Таблица~\ref{tab:constructs} связывает основные величины с наблюдением, на котором может основываться вывод. Такое разграничение необходимо уже при подготовке данных: метка воздействия, запись нейронного ответа и поведенческий исход задают разные задачи обучения.

\begingroup
\ResearchTableFont
\begin{longtable}{|L{37mm}|L{57mm}|L{65mm}|}
\caption{Наблюдаемые величины и границы их интерпретации}\label{tab:constructs}\\
\hline
Величина & Наблюдение & Обоснованный вывод \\
\hline
\endfirsthead
\multicolumn{3}{l}{Продолжение таблицы~\thetable}\\
\hline
Величина & Наблюдение & Обоснованный вывод \\
\hline
\endhead
\endfoot
\endlastfoot
Физическое воздействие & Температура, сила, длительность и место приложения & Условия входа; интенсивность стимула сама по себе не задаёт субъективную оценку \\
\hline
Нейронный ответ Drosophila & Ток, разряды, кальциевый показатель & Активность выбранных клеток с учётом способа регистрации \\
\hline
Защитное поведение & Перекатывание, сокращение, избегание, их длительность & Двигательная реакция в определённой экспериментальной пробе \\
\hline
Человеческий ECAP & Разность потенциалов после стимуляционного импульса & Вызванный суммарный ответ с учётом геометрии и артефактов \\
\hline
Клиническое состояние & Самоотчёт об интенсивности боли, функция, качество жизни & Исход в заданной популяции и в определённый срок наблюдения \\
\hline
\end{longtable}
\endgroup

\subsection{Нейронные цепи, защитная реакция и границы межвидового вывода}

Ноцицептивная реакция зависит не только от периферического обнаружения стимула, но и от состояния центральной сети. В исследовании острой ноцицепции и гиперчувствительности после повреждения показано участие разных нейронов и молекулярных механизмов. Поэтому изменение порога или длительности поведения нельзя заранее объяснять единственным общим параметром чувствительности. Для модели важно различать исходное состояние, воздействие и изменения, связанные с предшествующим повреждением \cite{s901}.

Нисходящие связи могут регулировать завершение защитного движения и переход к следующему действию. В исследовании GABA-ергических нейронов личинки анализировалось, как тормозная обратная связь участвует в прекращении перекатывания и изменении задержки последующего ползания. Результат касается организации последовательности поведения. Он обосновывает включение обратных связей в рассмотрение динамики, но не задаёт универсальный коэффициент торможения для любой ноцицептивной сети \cite{s910}.

Данные по взрослой мухе образуют самостоятельное основание. В препринте о восходящих ноцицептивных путях взрослой \textit{Drosophila} сопоставлены сенсорные реакции, оптогенетические вмешательства, поведение и анатомические связи. В литературном анализе используется закреплённая первая версия этой работы. Её выводы относятся к взрослой стадии; принадлежность к одному виду не делает найденные пути готовой заменой личиночным цепям. Возраст, строение тела, набор доступных нейронов и способ воздействия должны учитываться при выборе исходной модели \cite{s029}.

Отдельного внимания требует согласование единицы наблюдения. Несколько зарегистрированных клеток одного животного позволяют изучать пространственную организацию ответа, а повторные пробы --- его изменчивость в пределах этого животного. Эти данные не дают автоматически такого же основания для обобщения, как наблюдения у независимых животных. Аналогично число кадров видеозаписи не равно числу независимых поведенческих испытаний. Извлечение обучающих окон должно сохранять сведения об их принадлежности животному, сеансу и стимулу \cite{s212,s910}.

В контексте диссертации ноцицепция рассматривается через измеримые нейронные и защитные реакции. Самоотчёт человека об интенсивности боли является иной наблюдаемой величиной. В клинических исследованиях стимуляции спинного мозга изменения боли, функции и качества жизни регистрируются отдельно от электрического ответа нервных волокон. Такая организация измерений позволяет обсуждать связь технического воздействия с состоянием пациента без отождествления этих показателей \cite{s780,s781}.

Следовательно, исходная задача моделирования Drosophila должна иметь собственный проверяемый результат: воспроизведение выбранной характеристики активности или реакции на воздействие. Предположение о полезности полученных параметров для человеческой электрофизиологической модели формулируется на следующем уровне анализа. Для него требуется самостоятельное сравнение; оно не следует из наличия у мухи ноцицептивной цепи.

\section{Коннектомы Drosophila и анатомические ограничения нейродинамических моделей}
\label{sec:connectomes}

\subsection{Область охвата коннектома, версия и происхождение графа}

Коннектом задаёт сведения о нейронах и связях между ними в определённом анатомическом материале. При его использовании необходимо различать изображение ткани, реконструированные клетки, обнаруженные синапсы и производную матрицу для расчётов. Каждый переход предполагает отбор и обработку данных: объединение фрагментов, контроль реконструкции, выбор типов связей, порогов и области нервной системы. Поэтому название ресурса недостаточно для воспроизведения входа модели \cite{s904,s774,s738}.

В публикации о коннектоме мозга личинки описан граф с 3016 нейронами и приблизительно 548 тысячами синапсов. Это анатомическое описание, которое позволяет анализировать маршруты взаимодействий и организацию сети. Число синапсов между двумя клетками может использоваться как структурный признак связи, однако переход от этого числа к физической эффективности передачи требует дополнительных предположений. Сам анатомический результат не определяет кинетику синапсов или реакцию мозга на произвольное воздействие \cite{s904}.

Модели взрослого мозга используют иной анатомический материал. В вычислительной работе Shiu и соавторов сведения о связности и типах нейромедиаторов объединены в модели интегрирования с утечкой и пороговым формированием импульсов. Проверялись выбранные сенсомоторные преобразования. Значение этой работы состоит в демонстрации того, как крупный структурный ресурс может быть включён в функциональную модель при явно заданных правилах динамики \cite{s775}.

MANC охватывает вентральную нервную цепочку взрослого самца, а не весь мозг личинки. В настоящем обзоре используется первая версия рецензируемого препринта eLife. Наличие общедоступного анатомического ресурса ещё не устанавливает, какой именно экспорт был преобразован в матрицу конкретной программной реализации. Идентификатор набора, правила фильтрации и соответствие индексов узлам должны проверяться по входным файлам модели \cite{s774}.

В работе по BANC рассматривается анатомия мозга и вентральной нервной цепочки взрослой самки; в проверенном описании используется версия v888. Она расширяет область охвата по сравнению с ресурсом только мозга или только вентральной цепочки. При этом анализ влияний по графу характеризует организацию связей и не является сам по себе калиброванной биофизической моделью. Сопоставление предварительной и журнальной версий требует проверки состава данных, поскольку изменение числа узлов и трактовки связей может менять вход последующих расчётов \cite{s738}.

Таблица~\ref{tab:connectomes} показывает, почему эти ресурсы следует выбирать по исследовательской задаче, а не по одному общему признаку полноты. Полный граф выбранной области и полная нервная система организма --- разные характеристики охвата.

\begingroup
\ResearchTableFont
\begin{longtable}{|L{30mm}|L{42mm}|L{42mm}|L{39mm}|}
\caption{Анатомические ресурсы и условия их использования в модели}\label{tab:connectomes}\\
\hline
Ресурс или работа & Биологический материал & Назначение & Условие сопоставления \\
\hline
\endfirsthead
\multicolumn{4}{l}{Продолжение таблицы~\thetable}\\
\hline
Ресурс или работа & Биологический материал & Назначение & Условие сопоставления \\
\hline
\endhead
\endfoot
\endlastfoot
Мозг личинки \cite{s904} & Личиночная стадия; реконструкция мозга & Структурные ограничения нейронных путей & Возраст и наблюдения согласуются с выбранной моделью \\
\hline
Модель взрослого мозга \cite{s775} & Взрослая Drosophila; граф мозга & Сенсомоторное вычислительное моделирование & Фиксируются исходный граф и правила нейронной динамики \\
\hline
MANC \cite{s774} & Вентральная нервная цепочка взрослого самца & Анатомия двигательных цепей & Требуется точная версия экспорта и перечень выбранных узлов \\
\hline
BANC \cite{s738} & Мозг и вентральная нервная цепочка взрослой самки & Связи между отделами нервной системы & Журнальная версия и производная матрица проверяются раздельно \\
\hline
Совмещённые наблюдения \cite{s212} & Первичная личиночная реконструкция и функциональные пробы & Сопоставление анатомии с кальциевыми ответами & Возрастные группы и экспериментальные показатели сохраняются раздельно \\
\hline
\end{longtable}
\endgroup

Производные представления графа могут содержать другое число узлов, чем исходная публикация. Например, в препринте о замороженном частотном операторе рассматривается матрица на 3013 узлах, тогда как первичный анатомический источник сообщает 3016 нейронов. Эти величины относятся к разным описаниям данных. Для воспроизведения расчёта нужен перечень сохранённых идентификаторов и документированные преобразования, а не предположение о тождестве матрицы исходной реконструкции \cite{s904,s854}.

\subsection{Структурные ограничения и неопределённость динамики}

Общее описание динамической модели удобно представить в виде
\begin{equation}
 x_{t+\Delta t}=F_{\theta}(x_t,u_t;G), \qquad
 o_t=Q_{\psi}(x_{\leq t})+\varepsilon_t,
 \label{eq:state-observation}
\end{equation}
где $G$ --- выбранный граф; $x_t$ --- состояние модели; $u_t$ --- воздействие; $\theta$ --- параметры динамики; $Q_{\psi}$ --- оператор наблюдения с параметрами $\psi$; $o_t$ --- измеряемая величина; $\varepsilon_t$ --- ошибка наблюдения. Формула~\eqref{eq:state-observation} разделяет уравнение состояния и способ получения экспериментального показателя. Для кальциевой регистрации $Q_{\psi}$ должен учитывать особенности индикатора и временной фильтрации, а для электрической регистрации --- соответствующий измерительный канал \cite{s212,s285}.

Задание $G$ сужает множество возможных моделей, но оставляет неопределёнными постоянные времени, усиления, пороги, внешние входы и другие параметры. Даже знак влияния связи может требовать сведений о нейромедиаторе и рецепторе. Если модель использует только положительные числа синапсов, полученная матрица описывает иной оператор, чем сеть с установленными возбуждающими и тормозными взаимодействиями. Сравнивать их следует с учётом этих различий, а не только совпадения топологии \cite{s743,s854}.

В модели FlyVis анатомические ограничения сочетаются с обучением неизвестных параметров и проверкой зрительных ответов. Результат показывает продуктивность такого сочетания для определённой системы и задачи. Он не устанавливает, что любой коннектом при произвольной параметризации воспроизводит активность живой нервной системы. Область наблюдений и обучающая задача остаются частью модели \cite{s744}.

Теоретическая работа Beiran и Litwin-Kumar уточняет условия восстановления активности в коннектомно-ограниченных рекуррентных сетях. У обучаемой сети и вычислительного «учителя» заданы общие синаптические веса, входы и начальные состояния. При таких предпосылках исследуется, что можно восстановить по наблюдениям части нейронов. Вычислительный эталон позволяет проверить определимость в заданном классе систем, но его нельзя считать экспериментальной записью, подтверждающей биологическую динамику \cite{s743}.

Иной тип сопоставления представлен в работе, объединяющей кальциевое изображение и электронную микроскопию личиночной нервной системы. Первичные совместные данные относятся к первому личиночному возрасту; дополнительные функциональные и поведенческие проверки используют другие личинки третьего возраста. Это позволяет сопоставлять разные основания вывода, сохраняя различие выборок. Ответ выбранных нейронов на оптогенетическое воздействие и поведенческая реакция другой возрастной группы не образуют единого синхронного временного ряда \cite{s212}.

Наконец, вывод о роли структуры зависит от устройства контрольного сравнения. В препринте о чувствительности зрительной сети к топологии отдельно проверены влияние общей инициализации и перестройки графа с сохранением степеней. Изменение результата при переходе между контролями показывает, что преимущество исходного графа может смешиваться с преимуществом исходных параметров. Такой результат относится к проверенной задаче и не является универсальным отрицанием значения коннектома \cite{s741}.

Для диссертационного исследования отсюда следует требование раздельно фиксировать граф, параметры динамики и способ наблюдения. Их происхождение, область применимости и проверка должны быть доступны независимо друг от друга. Только после этого синтетические ряды можно использовать как воспроизводимый источник предобучения.

\section{Формирование ECAP и измерительный канал системы стимуляции спинного мозга}
\label{sec:ecap-channel}

\subsection{Возбуждение волокон и регистрация суммарного ответа}

Вызванный составной потенциал действия (evoked compound action potential, ECAP) представляет собой регистрируемый электрический ответ совокупности активированных нервных волокон. При стимуляции спинного мозга (spinal cord stimulation, SCS) на его формирование влияют распределение внешнего поля, свойства волокон, распространение потенциалов действия и конфигурация отведения. Поэтому задача расчёта сигнала включает несколько связанных физических процессов \cite{ecapmodel,s285}.

На первом этапе определяется поле, создаваемое электродами в объёмном проводнике. Расчёт зависит от геометрии тканей и контактов, электрической проводимости, формы импульса и граничных условий. Далее это поле используется как внешнее воздействие на модель волокна. Появление и распространение потенциала действия задаются свойствами мембраны и геометрией аксона. Наконец, токи активированных волокон формируют потенциалы на регистрирующих контактах. Вычисление возбуждения и вычисление измеренного напряжения имеют разные входы и должны проверяться по соответствующим наблюдениям \cite{ecapmodel,posture}.

При линейном описании пассивного объёмного проводника связь с регистрацией можно схематически представить суммой вкладов сегментов волокон:
\begin{equation}
 \widehat y_{\mathrm n}(t;\gamma)=
 \sum_{q=1}^{N_{\mathrm a}}\sum_{k=1}^{M_q}
 Z_{qk}(\gamma) I_{m,qk}(t),
 \label{eq:recording-operator}
\end{equation}
где $I_{m,qk}$ --- трансмембранный ток сегмента $k$ волокна $q$; $Z_{qk}$ --- коэффициент передачи к выбранному отведению; $\gamma$ --- параметры геометрии, проводимости и регистрации; $N_{\mathrm a}$ --- число моделируемых волокон; $M_q$ --- число сегментов волокна. Если ток задан в амперах, а коэффициент передачи --- в омах, результат имеет размерность напряжения. Формула~\eqref{eq:recording-operator} отражает принцип суммирования, используемый в прямых и основанных на взаимности моделях регистрации, при принятых допущениях о проводнике \cite{ecapmodel,posture}.

Из такого представления следует, что число возбуждённых волокон не определяет сигнал однозначно. Имеют значение их положение относительно электродов, диаметры, распределение токов и времена прихода ответов. Две совокупности активированных волокон могут давать близкую амплитуду при различном составе; сходное рекрутирование может сопровождаться разными напряжениями при изменении измерительной геометрии. Эту неоднозначность непосредственно рассматривает исследование влияния позы и ширины импульса на ECAP \cite{posture}.

В реальной записи к нейронной составляющей добавляются стимуляционный артефакт, аппаратный шум и другие компоненты. Поэтому наблюдение следует описывать как
\begin{equation}
 y(t)=y_{\mathrm n}(t)+a(t)+\varepsilon(t),
 \label{eq:raw-recording}
\end{equation}
где $a(t)$ обозначает артефакт, а $\varepsilon(t)$ --- совокупную ошибку регистрации. Разделение в формуле~\eqref{eq:raw-recording} является модельным: по одной форме волны происхождение каждого вклада может быть не установлено. В техническом исследовании на свиньях показаны ненейронные компоненты, способные имитировать признаки вызванного ответа. Этот результат обосновывает необходимость контрольных проверок регистрации; человеческие клинические выводы из такой модели не следуют \cite{s798}.

Рисунок~\ref{fig:ecap-formation} связывает перечисленные этапы. Обработка относится к уже зарегистрированной смеси компонентов. Она может облегчать выделение ответа, но её качество оценивается отдельно от корректности расчёта возбуждения и от эффективности терапии.

\begin{figure}[htbp]
\centering
% Literature synthesis: separate excitation, recording and signal processing.
\begin{tikzpicture}[x=1mm,y=1mm]
\node[research process] (pulse) at (0,0) {Параметры\\импульса};
\node[research process] (field) at (55,0) {Внешнее\\электрическое поле};
\node[research process] (axon) at (110,0) {Возбуждение\\волокон};
\node[research process] (curve) at (0,-34) {Кривая роста\\ECAP};
\node[research process] (process) at (55,-34) {Обработка\\записи};
\node[research process] (recording) at (110,-34) {Оператор\\регистрации};
\node[draw,circle,inner sep=0pt,minimum size=3.4mm,line width=0.5pt,
  font=\fontsize{13}{15.7}\selectfont] (sum) at (82.5,-34) {$+$};
\node[font=\fontsize{13}{15.7}\selectfont,align=center,inner sep=0pt]
  (noise) at (82.5,-17) {Артефакт и шум};
\draw[research flow] (pulse) -- (field);
\draw[research flow] (field) -- (axon);
\draw[research flow] (axon) -- (recording);
\draw[research flow] (recording) -- (sum);
\draw[research flow] (sum) -- (process);
\draw[research flow] (noise) -- (sum);
\draw[research flow] (process) -- (curve);
\end{tikzpicture}

\caption{Последовательность формирования и обработки ECAP}
\label{fig:ecap-formation}
\end{figure}
\FloatBarrier

\subsection{Кривая роста, геометрия и стимуляционный артефакт}

Основной целевой величиной в постановке диссертации является кривая роста --- зависимость амплитуды ECAP от интенсивности стимуляции. Её определение требует фиксированных условий: формы и ширины импульса, частоты, конфигурации контактов, позы, правил подавления артефакта и выделения амплитуды. Изменение этих условий означает изменение измерительной задачи, даже если график по-прежнему строится в координатах «ток --- напряжение» \cite{posture,s787}.

В общем виде амплитудную характеристику можно записать как
\begin{equation}
 A(I;\gamma)=\mathcal M_{\mathcal W}
 \bigl[y_{\mathrm n}(t;I,\gamma)\bigr],
 \label{eq:growth-curve}
\end{equation}
где $I$ --- амплитуда стимуляционного тока; $\mathcal W$ --- окно анализа; $\mathcal M_{\mathcal W}$ --- принятый способ измерения амплитуды. В зависимости от протокола он может использовать разность выбранных пиков либо другой показатель выделенного ответа. Формула~\eqref{eq:growth-curve} показывает, что кривая зависит не только от возбуждения, но и от способа получения скалярного показателя. При сравнении моделей этот способ должен быть одинаковым \cite{posture,s159}.

Порог и начальный наклон кривой дают компактные характеристики ответа в определённом диапазоне воздействия. Локальная кусочно-линейная аппроксимация является естественной моделью сравнения: она описывает изменение амплитуды без скрытого временного состояния. Её область пригодности ограничена диапазоном данных; насыщение, нестабильное выделение пиков и изменение параметров импульса могут нарушать простую зависимость. Усложнение модели оправдывается сравнением с таким исходным описанием, а не предполагаемой выразительностью нейронной сети \cite{s159,posture}.

Измерительные условия могут существенно менять связь между кривой роста и рекрутированием. В исследовании позы и ширины импульса первоначально участвовали 56 человек, а для анализа кривых после исключений использованы данные 45 участников. Применение вычислительной модели позволило отдельно рассмотреть влияние геометрии на записанный потенциал и на активацию волокон. В результате постоянство амплитуды ECAP не могло использоваться как достаточное условие постоянства рекрутирования во всех рассмотренных конфигурациях \cite{posture}.

Для обработки важен компромисс между уменьшением артефакта и сохранением нейронного ответа. Модель артефакта, построенная на аддитивности и граничных условиях, исследовалась на большом числе сигналов от двух пациентов. Многочисленные записи позволяют изучать повторяемость технической процедуры, но не увеличивают число независимых участников. Кроме того, в доступном тексте работы определения итогового E-показателя в основном изложении и приложении расходятся. Его численные значения поэтому не используются здесь для ранжирования методов \cite{s159}.

Сходное разграничение необходимо для исследований автоматизированного программирования SCS. Работа о следующем поколении программирующей платформы оценивает технические показатели и пользовательский опыт; исходы эффективности терапии в рассмотренных исследованиях не собирались. В сравнении с ней исследования EVOKE регистрируют клинические и функциональные исходы. Эти результаты отвечают на разные вопросы и требуют разных эталонов \cite{s162,s780,s781}.

Для моделирования кривой роста недостаточно подобрать параметры так, чтобы расчётная линия совпала с записью одного участника. Следует проверить устойчивость описания при изменении условий, установить вклад регистрации и заранее определить, какие исходные сведения будут доступны для нового случая. Технический результат при этом выражается ошибкой воспроизведения электрического показателя; вывод об изменении интенсивности боли требует самостоятельных связанных наблюдений.

\section{Методы моделирования нейродинамики, формирования ECAP и переноса представлений}
\label{sec:models-transfer}

\subsection{Биофизические, частотные и обучаемые модели}

Выбор класса модели определяется тем, какая величина должна быть воспроизведена и по каким данным будут определяться параметры. Многокомпартментные модели аксонов описывают возбуждение и распространение потенциалов действия; объёмный проводник связывает эти процессы с электродным напряжением. Частотные и импульсные сетевые модели используют другое разрешение описания. Обучаемая модель может аппроксимировать отдельный расчётный блок, но её выход и область проверки должны оставаться определёнными \cite{ecapmodel,s285}.

Коннектомное ограничение совместимо с разными уравнениями состояния. В модели интегрирования с утечкой учитываются накопление входного воздействия, спад состояния, достижение порога и формирование импульса. Частотная модель вместо отдельных разрядов описывает выбранную усреднённую активность. Первая ближе к последовательностям событий, вторая может облегчать исследование крупных сетей. Физическая интерпретация параметров в обоих случаях зависит от принятого масштаба времени и способа согласования с наблюдениями \cite{s775,s212}.

В работе о кольцевом аттракторе анатомические сведения используются для построения динамики ориентационной системы. В препринте о ходьбе анализируются свойства вычислительной двигательной сети. Эти примеры показывают, что функциональную постановку необходимо задавать дополнительно к графу: устойчивое представление направления, последовательность активности двигательных узлов и ноцицептивная реакция --- разные выходы. Предварительный статус работ и ограничения функциональных наблюдений сохраняются при их сопоставлении \cite{s900,s909}.

Для моделей с телом появляется дополнительная связь между нейронным выходом, механикой движения и сенсорной обратной связью. В препринте о коннектомном графовом контроллере локомоция получена в симулированной среде с использованием обучения с подкреплением. Такой результат характеризует работу вычислительного контроллера и тела. Его биологическая достоверность требует сопоставления с независимыми наблюдениями животного, а перенос к человеческому измерительному каналу --- ещё одного основания \cite{s286}.

Таким образом, будущий генератор нейродинамических данных должен различать состояние сети, состояние тела и наблюдение среды. Наличие анимации движения не подтверждает достоверность внутренней нейронной активности. С другой стороны, воспроизведение сетевой активности без тела не проверяет замкнутый сенсомоторный механизм. Выбранный уровень моделирования должен соответствовать исходной обучающей задаче; остальные компоненты могут оставаться за её пределами.

Ускорение человеческой модели также возможно без обращения к другому биологическому виду. В работе об обобщённой активационной функции исследуется приближённый способ оценки возбуждающего действия поля. Сопоставление проводится с многокомпартментными симуляциями. В проверенной принятой рукописи описаны ограничения подпорогового режима: такая функция не охватывает все активные мембранные эффекты, рефрактерность и блок проведения. Поэтому вычислительная точность относительно выбранного симулятора ещё не означает пригодность для всех временных протоколов или наличие независимо проверенного ECAP \cite{s877}.

Результаты этих работ целесообразно сопоставлять по двум признакам: какой процесс моделируется непосредственно и какой показатель служит эталоном. Проверка реализации по другой программе относится к вычислительному уровню. Сравнение с нейронной записью проверяет соответствие конкретному наблюдению. Оценка на новых участниках касается обобщения. Эти основания дополняют друг друга, но не образуют единого показателя качества \cite{s285,s743}.

\subsection{Синтетические данные и перенос начальных параметров}

Синтетические данные могут использоваться для изучения свойств модели, обучения приближённого вычислителя и предварительной настройки параметров. Их преимущество состоит в управляемости условий генерации и доступности внутренних состояний. Одновременно они наследуют допущения симулятора. Если целевая выборка отличается от генератора по геометрии, шуму или механизму формирования метки, увеличение объёма синтетических наблюдений само по себе не устраняет это расхождение \cite{s285,s836}.

В исследовании синтетического дополнения электроэнцефалографических данных показана необходимость контролировать изменение качества при переходе к новым участникам. Полезность синтетического примера рассматривается относительно целевой задачи и допустимых условий его отбора. Данный результат относится к использованным наборам и признакам; различение клинической группы и контроля в таком наборе не тождественно оценке интенсивности боли, ответа на SCS или амплитуды ECAP \cite{s836}.

Работа о синтетическом обучении мультимодальной оценке боли представляет иной режим: реальные неразмеченные данные используются для адаптации. В описанном протоколе общий пул применялся во всех разбиениях без строгого исключения данных тестовых участников. Авторы отмечают возможное оптимистическое смещение, но его размер не установлен. Из этого следует ограничение независимости оценки; фактическое пересечение участников нельзя считать измеренным только на основании устройства протокола \cite{s149}.

Для переноса требуется определить, что сохраняется между исходной и целевой задачами. В работе о представлении коннектомов для резервуарных вычислений латентное описание графов связано с качеством вычислительных задач. Использован ресурс коркового коннектома мыши. Это пример обучения представления структуры и оценки его функционального значения; он отличается от предобучения временного модуля на наблюдениях мухи и от расчёта человеческого электрического ответа \cite{s892}.

В доступной журнальной аннотации работы о переносе между мышью и человеком рассматривается соответствие электрофизиологических признаков и транскриптомных классов корковых интернейронов. По этому материалу можно установить область задачи, но нельзя восстановить все условия настройки и проверки. Этот пример показывает, насколько существенно указание целевого выхода: перенос для классификации клеток и перенос для моделирования составного потенциала требуют разных данных и критериев \cite{s873}.

В постановке настоящего исследования переносимым объектом выступают начальные параметры временного модуля. Анатомический граф мухи, её клетки и внутренние состояния не получают интерпретации человеческих волокон. Исходный модуль может обучаться прогнозу последующей активности по доступной истории наблюдений и воздействий. Будущий целевой отсчёт в такую историю не включается. В целевой модели используются собственные входное преобразование, описание возбуждения и оператор регистрации. Это определяет проверяемую гипотезу сопряжения, а не свидетельство её полезности.

Необходимость временного представления зависит от состава целевых данных. Кривая, построенная только по амплитуде изолированного импульса при фиксированных условиях, может описываться без памяти. История стимуляции, интервалы между импульсами и восстановление возбудимости задают возможные временные зависимости. Их наличие должно быть установлено в пригодных записях; вводить рекуррентный модуль исключительно потому, что исходная модель мухи динамическая, методически недостаточно \cite{posture,s877}.

Отсюда следуют две самостоятельные проверки. Первая устанавливает, нужно ли временное состояние для целевой задачи по сравнению с простой кривой роста. Вторая выясняет, улучшает ли исходное предобучение работу того же целевого класса моделей. При таком разделении можно отличить эффект выбранной архитектуры от эффекта начальных параметров и конкретного источника динамики.

\section{Сравнительный анализ существующих подходов и выбор методов сравнения}
\label{sec:comparisons}

\subsection{Ближайшие методические аналоги и область установленного результата}

Сопоставление рассмотренных работ выполняется по задаче, исходным данным и виду проверки. Для анатомического ресурса существенны область реконструкции и версия. Для динамической модели --- доступные функциональные наблюдения и контролируемые альтернативы. Для расчёта ECAP --- описание возбуждения, регистрации и геометрии. Для переноса --- данные предобучения, адаптации и независимой оценки. Такое направленное сопоставление позволяет проследить основания конкретного вывода; полнота систематического охвата всей литературы здесь не предполагается.

Таблица~\ref{tab:analogues} объединяет ближайшие подходы к отдельным частям постановки. Журнальные публикации, препринты и патентная заявка сохраняют свой статус. Преимущество, установленное в конкретной зрительной или вычислительной задаче, относится к соответствующему столбцу таблицы, а не к будущей задаче моделирования ECAP.

\begingroup
\ResearchTableFont
\begin{longtable}{|L{39mm}|L{39mm}|L{40mm}|L{35mm}|}
\caption{Методические аналоги и область их проверки}\label{tab:analogues}\\
\hline
Подход & Данные или эталон & Установленный результат & Значение для постановки \\
\hline
\endfirsthead
\multicolumn{4}{l}{Продолжение таблицы~\thetable}\\
\hline
Подход & Данные или эталон & Установленный результат & Значение для постановки \\
\hline
\endhead
\endfoot
\endlastfoot
Модель взрослого мозга \cite{s775} & Коннектом и выбранные сенсомоторные наблюдения & Функциональные преобразования при заданных правилах динамики & Пример исходной сети; параметры и область проверки фиксируются \\
\hline
FlyVis \cite{s744} & Зрительная связность и наблюдаемые ответы & Обучение неизвестных параметров при анатомических ограничениях & Пример совместного использования структуры и функции \\
\hline
Личиночная модель \cite{s212} & Анатомия и кальциевые показатели & Сопоставление знаков и характера ответов & Требуется согласовать возраст и оператор наблюдения \\
\hline
Рекуррентная сеть «учитель --- ученик» \cite{s743} & Вычислительный эталон с общими весами, входами и начальными состояниями & Определимость активности и параметров при заданных предпосылках & Основание для анализа идентифицируемости \\
\hline
Контроли топологии \cite{s741,s854} & Зрительная задача и замороженный личиночный оператор; препринты & Зависимость вывода от инициализации, степеней и весов & Раздельная оценка структуры и обученных параметров \\
\hline
Прямое моделирование ECAP \cite{ecapmodel,posture} & Поле, аксоны, запись при заданной геометрии & Связь стимуляции, возбуждения и измеренного ответа & Основной человеческий физический компонент \\
\hline
Обобщённая активационная функция \cite{s877} & Сопоставление с многокомпартментным симулятором & Ускоренная оценка возбуждения в проверенном режиме & Человеческая вычислительная альтернатива; регистрация проверяется отдельно \\
\hline
Синтетическое обучение и адаптация \cite{s149} & Синтетический источник и неразмеченный реальный пул & Качество в описанном протоколе оценки боли & Независимость от тестовых участников не обеспечена \\
\hline
Латентное представление коннектомов \cite{s892} & Корковые графы мыши и резервуарные задачи; препринт & Связь представления структуры с вычислительным качеством & Пример другого переносимого объекта и другой задачи \\
\hline
Синтетические ECAP и предобучение \cite{s746} & Патентное описание и формула способа & Раскрыта схема обучения классификатора & Предшествующее техническое решение; эффективность отдельно не установлена \\
\hline
\end{longtable}
\endgroup

В совокупности работы показывают наличие отдельных необходимых компонентов: коннектомных ограничений, функциональной настройки, генерации синтетических наблюдений и физических моделей регистрации. При этом одинаковый термин «нейронная модель» обозначает разные объекты. Сеть, воспроизводящая сенсомоторный выбор, и аксонная модель, определяющая токи для электрического отведения, не имеют общего выходного пространства без дополнительного преобразования. Именно это различие определяет задачу сопряжения \cite{s744,s775,ecapmodel}.

Патентная заявка WO2025224687A1 описывает синтетические ECAP и предобучение классификатора с использованием схемы кодирования и восстановления. Поэтому само сочетание синтетических сигналов и предобучения уже имеет технические аналоги. Патентное описание устанавливает факт раскрытия способа; оно не предоставляет независимой участнической оценки, которая могла бы подтвердить преимущество исследуемого переноса \cite{s746}.

Предполагаемый вклад диссертации следует связывать с конкретным способом задания начальных параметров временного модуля, его сопряжением с человеческой физической моделью и условиями дополнительной точности. Перечисление известных блоков либо отсутствие найденной работы с тем же названием не устанавливает научную новизну. Ограничения рассмотренного корпуса сохраняют значение и после выбора ближайших аналогов: вопрос приоритета и вопрос полезности конкретного метода требуют самостоятельных оснований.

\subsection{Контрольные сравнения и разделение возможных эффектов}

Сравнение должно отвечать на определённый вопрос об эффекте. Если у двух вариантов различаются архитектура, состав входов и число шагов настройки, итоговая разность ошибок характеризует всё это сочетание. Для оценки предобучения необходимо сохранить сопоставимый целевой модуль, одинаковую доступную информацию и общий протокол выбора параметров. Проверка роли графа дополнительно требует согласованного способа инициализации и описания свойств случайного контроля \cite{s741,s743}.

Таблица~\ref{tab:review-controls} задаёт состав сравнений, вытекающий из литературного анализа и цели исследования. Она определяет проверяемые эффекты, но не подменяет выбор окончательного протокола после подтверждения пригодности данных.

\begingroup
\ResearchTableFont
\begin{longtable}{|L{48mm}|L{52mm}|L{59mm}|}
\caption{Контрольные сравнения для оценки предполагаемого переноса}\label{tab:review-controls}\\
\hline
Вариант & Проверяемый эффект & Условие интерпретации \\
\hline
\endfirsthead
\multicolumn{3}{l}{Продолжение таблицы~\thetable}\\
\hline
Вариант & Проверяемый эффект & Условие интерпретации \\
\hline
\endhead
\endfoot
\endlastfoot
Простая кривая роста без скрытого временного состояния & Необходимость временного представления & Одинаковые физические входы; использование истории оговаривается явно \\
\hline
Целевая модель без предобучения & Вклад начальных параметров & Сопоставимые размерность, настройка и доступные данные человека \\
\hline
Предобучение на ECAP-подобных сигналах человеческой модели & Выбор исходного домена & Одинаковые правила генерации, объём данных и целевая проверка \\
\hline
Неноцицептивная динамика Drosophila & Специфичность исходной реакции & Сопоставимые обучающая задача, объём рядов и условия настройки \\
\hline
Перемешивание временной организации & Значение порядка наблюдений & Сохранение распределений и явно заданное правило перестановки \\
\hline
Перестроенные коннектомы & Значение топологии & Контроль степеней, весов, знаков и инициализации в пределах выбранной модели \\
\hline
Структура без обученных динамических параметров & Разделение топологии и обучения & Независимое описание того, какие параметры заданы и какие оценивались \\
\hline
Замороженный и дообучаемый временной модуль & Необходимость целевой адаптации & Общие разбиения и сопоставимые условия настройки остальных компонентов \\
\hline
Вариант без оператора регистрации & Возможность нефизического получения низкой ошибки & Только диагностическое сравнение; не подтверждает биофизический механизм \\
\hline
\end{longtable}
\endgroup

Сохранение числа рёбер само по себе недостаточно для графового контроля. При перестройке могут меняться распределения входных и выходных степеней, суммарные веса, соотношение возбуждения и торможения, спектральные свойства и масштаб ответа. Препринт о замороженном операторе отдельно показывает зависимость различных показателей от разных ограничений ансамбля. Поэтому выражение «случайный граф» должно сопровождаться точным определением сохранённых характеристик \cite{s854}.

Аналогично перестановка времени не является единственным универсальным контролем динамики. Перемешивание отдельных каналов, совместных отсчётов или целых окон по-разному меняет доступную мгновенную информацию и прогнозную задачу. Для содержательного сравнения нужно заранее указать операцию и проверить, какие свойства данных она сохраняет. Иначе разность результатов будет включать изменение самой обучающей задачи.

Контроль без физического оператора имеет особое назначение. Он может показать, позволяет ли гибкая регрессия получать малую ошибку без соблюдения выбранного механизма регистрации. Такая возможность помогает обнаружить недостаточность проверочных данных или зависимость результата от нефизического преобразования. Для подтверждения метода биофизического моделирования требуется дополнительное соответствие физическим ограничениям, а не только преимущество по ошибке.

\section{Требования к исходным данным и экспериментальной проверке}
\label{sec:data-validation}

\subsection{Происхождение данных, совместимость и единица наблюдения}

Требования к данным определяются совместно исходной и целевой задачами. Для Drosophila необходимы версия графа, перечень идентификаторов, правила фильтрации и сведения о биологическом материале. Функциональные наблюдения должны описывать возраст, воздействие, регистрируемые клетки и способ измерения. Различие между состоянием модели и наблюдением, введённое в формуле~\eqref{eq:state-observation}, должно сохраняться и в обучающих данных \cite{s904,s212}.

Генератор синтетической динамики должен позволять связать каждый ряд с графом, параметрами, воздействием и начальным состоянием. Для воспроизведения существенны также шаг интегрирования, случайные начальные значения, моменты сброса и способ формирования наблюдаемого выхода. Временные ряды сети, координаты тела и параметры среды могут храниться совместно, но иметь разные физические единицы и частоты дискретизации. Их синхронизация является частью методики, а не следствием общего имени файла.

Из сопоставления нейродинамических и электрофизиологических моделей следует различие двух видов синтетических данных. Первый представляет активность исходной сети и служит для обучения временного модуля. Второй представляет целевой сигнал, рассчитанный при заданных характеристиках человеческих волокон и регистрации. Объединение этих видов данных в одном вычислительном процессе не устраняет различия их физического смысла. Поэтому происхождение записи должно позволять восстановить не только случайное начальное значение, но и тот оператор, который преобразовал состояние модели в наблюдение \cite{s285,ecapmodel,s746}.

Для оценки устойчивости необходимо различать повторение одного расчёта и изменение условий модели. Повтор с теми же входами проверяет вычислительную воспроизводимость в принятом численном режиме. Изменение структуры, параметров, воздействий и начальных состояний позволяет исследовать зависимость результата от выбранной модели. Даже большая совокупность таких расчётов остаётся выборкой из заданного семейства симуляций. Она не устанавливает автоматически соответствие распределению биологических наблюдений; это соответствие требует отдельного сопоставления с измерениями. Смешение этих уровней могло бы привести к тому, что оценка переноса характеризовала бы особенности общего симулятора вместо полезности исходного представления \cite{s743,s836}.

Для человеческого ECAP требуется идентифицировать участника, сеанс, конфигурацию электродов, параметры стимуляции и единицы напряжения и времени. Кривая роста должна сопровождаться исходными настройками и правилом выделения амплитуды. При использовании обработанных признаков нужно знать, какие процедуры уже применены к сигналу и можно ли повторить их независимо. Особенно важно установить, доступны ли действительные записи, а не только графики или агрегированные показатели публикации \cite{posture,s105}.

Таблица~\ref{tab:human-observations} показывает различие между числом записей и числом участников в основных рассмотренных работах. Она предназначена для оценки структуры доказательств, а не для выбора наибольшего массива по числу сигналов.

\begingroup
\ResearchTableFont
\begin{longtable}{|L{30mm}|L{45mm}|L{43mm}|L{35mm}|}
\caption{Человеческие наблюдения ECAP и границы независимости оценки}\label{tab:human-observations}\\
\hline
Работа & Участники и записи & Проверяемый вопрос & Граница обобщения \\
\hline
\endfirsthead
\multicolumn{4}{l}{Продолжение таблицы~\thetable}\\
\hline
Работа & Участники и записи & Проверяемый вопрос & Граница обобщения \\
\hline
\endhead
\endfoot
\endlastfoot
Поза и ширина импульса \cite{posture} & 56 участников первоначально; 45 в анализе кривых роста & Изменение сигнала и рекрутирования при разных условиях & Повторные воздействия внутри участников; не проверка переноса \\
\hline
Пороги восприятия \cite{s787} & 14 участников; 112 кривых роста & Связь ECAP с порогами восприятия и дискомфорта & Кривые не являются независимыми участниками; обезболивание не измеряется этим порогом \\
\hline
Признаки ECAP \cite{s105} & Из 15 привлечённых выбраны восемь; после исключений 28 сеансов & Классификация сигналов и анализ признаков & Внешнее разбиение по каждому одиннадцатому сигналу, а не по участнику \\
\hline
Выделение ответа \cite{s159} & Два участника; множество повторных записей & Удаление стимуляционного артефакта & Большое число сигналов не задаёт независимую межучастническую проверку \\*
\hline
\end{longtable}
\endgroup

Доступность данных устанавливается отдельно от доступности статьи. Для работ о позе и признаках ECAP первичные человеческие записи предоставляются по запросу с ограничениями; публикация не равнозначна открытому набору для загрузки. Подтверждение прав использования, состава записей и пригодности для независимой проверки должно предшествовать окончательному выбору модели. На этом основании для последующих этапов предусматривается выбор основного и независимого ECAP-наборов после проверки доступа \cite{posture,s105}.

Клиническая ветвь требует другого состава информации: связи между участником, электрофизиологией, параметрами терапии, показателями состояния и сроком наблюдения. Сообщения EVOKE о 24 и 36 месяцах относятся к одной исходной рандомизированной когорте. Рассматривать их как две независимые репликации было бы ошибкой. Способ обработки пропусков и последующий характер отдельных анализов также входят в интерпретацию результата \cite{s780,s781}.

\subsection{Независимая проверка, ошибка и неопределённость}

Верификация расчёта, настройка параметров и валидация по наблюдениям должны быть различимы в протоколе. Согласие двух реализаций уравнений поддерживает вычислительную корректность. Оценивание параметров по известным записям характеризует калибровку модели. Для вывода о применимости нужны наблюдения, не использованные при выборе этих параметров, и заранее определённая область проверки. Такая организация следует из модульного строения физической модели и из зависимости обучаемого результата от данных настройки \cite{s285,ecapmodel}.

Основной единицей статистического вывода для человеческой задачи является участник. Сигналы, окна, кривые и сеансы одного человека зависимы. Их распределение между обучающей и тестовой частями может давать хорошую оценку распознавания очередного сигнала в знакомых условиях, но не отвечает на вопрос о работе с новым участником. Пример классификации по каждому одиннадцатому сигналу показывает, почему название внешнего цикла проверки само по себе не устанавливает независимость на нужном уровне \cite{s105}.

Для основной оценки переноса разделение участников должно предшествовать выбору размерности, временного масштаба, нормировок и параметров обработки. Эти решения определяются внутри обучающих разбиений. Физические сведения, доступные при применении, например конфигурация электродов, могут входить в модель нового участника как заранее предусмотренные входы. Использование его зарегистрированных ответов для настройки представляет отдельный режим индивидуальной калибровки и требует отдельной оценки на оставшихся данных.

В качестве основного показателя будущего сравнения принимается ошибка амплитуды кривой роста в микровольтах. Для участника $i$ среднеквадратическая ошибка (root mean square error, RMSE) определяется по сопоставимым точкам кривой, затем ошибки усредняются между участниками:
\begin{equation}
 E_i=\sqrt{\frac{1}{m_i}\sum_{j=1}^{m_i}
       (\widehat A_{ij}-A_{ij})^2}, \qquad
 E=\frac{1}{N}\sum_{i=1}^{N}E_i.
 \label{eq:participant-error}
\end{equation}
Здесь $A_{ij}$ и $\widehat A_{ij}$ --- зарегистрированная и расчётная амплитуды; $m_i$ --- число сопоставимых точек; $N$ --- число участников. Формула~\eqref{eq:participant-error} задаёт одинаковый вес участникам, а не каждому зарегистрированному импульсу. При нескольких конфигурациях правило их объединения внутри участника должно быть определено заранее.

Калибровка модели по данным и калибровка прогнозных интервалов --- разные процедуры. Первая определяет параметры, вторая проверяет, соответствует ли фактическое покрытие интервала заявленному уровню. Интервалы, полученные в рамках одного симулятора, могут отражать только часть неопределённости. Изменение геометрии, состава участников или способа регистрации требует самостоятельной проверки устойчивости. Синтетическое дополнение и адаптация поэтому оцениваются на целевых данных в пределах явно установленного режима \cite{s285,s836,s149}.

Для последующего сравнения предлагается рассматривать разность ошибок контрольного метода и переноса. Положительная разность означает меньшую ошибку переноса. При заранее выбранном допуске практически значимого различия доверительный интервал позволяет различить четыре исхода. Если интервал целиком превышает положительный допуск, результат поддерживает практически значимое преимущество. Если он целиком находится внутри области от отрицательного до положительного допуска, различие признаётся практически несущественным в пределах принятой точности. Интервал целиком ниже отрицательного допуска указывает на ухудшение. Остальные положения интервала оставляют результат неопределённым относительно этих критериев. Эта схема является предложением для будущего протокола, а не результатом проведённого сравнения.

Отсутствие статистически значимого преимущества само по себе не устанавливает эквивалентность. Необходимые допуски, объём выборки и способ построения интервала должны быть определены после проверки пригодности данных и до обращения к окончательной тестовой части. Сравниваемые ошибки следует рассчитывать на одних и тех же участниках: иначе разность может отражать состав выборок. Равный вес участника в формуле~\eqref{eq:participant-error} также не устраняет различия числа точек и диапазонов стимуляции. Общий протокол должен фиксировать эти диапазоны, обработку отсутствующих точек и правило включения кривой в анализ. Участники с недостаточным числом пригодных измерений не могут незаметно исчезать из отчёта, поскольку исключения меняют область обобщения \cite{posture,s105}.

До независимой проверки разработанного оператора его вычислительный выход следует обозначать как ECAP-подобный сигнал. Отсутствие пригодных человеческих данных позволяет исследовать свойства компонентов и чувствительность расчётов, но ограничивает вывод об эффекте переноса. Техническая точность в воспроизведении кривой роста, её устойчивость и интервальная оценка составляют основную область будущей проверки. Клиническая эффективность и настройка терапии требуют дополнительного исследования на связанных исходах пациентов.

\section*{Выводы по главе}
\addcontentsline{toc}{section}{Выводы по главе}

Литературный анализ позволяет сформулировать следующие положения.
\begin{enumerate}
\item Физический стимул, активность нейронов, защитное поведение Drosophila, человеческий ECAP и клиническое состояние относятся к различным наблюдаемым величинам. Их разграничение определяет состав данных и допустимую интерпретацию результатов.
\item Коннектом ограничивает структуру модели, но оставляет неопределёнными динамические параметры и оператор наблюдения. Для исходной модели необходимы совместимые версия графа, стадия развития и функциональные данные; проверка соответствия выполняется отдельно от подготовки анатомической матрицы.
\item Расчёт ECAP требует согласованного описания внешнего поля, возбуждения, распространения и регистрации. Число активированных волокон не задаёт напряжение однозначно. Геометрия, артефакт и правило выделения амплитуды входят в постановку задачи моделирования кривой роста.
\item Рассмотренные методы подтверждают возможность использования структурных ограничений, обучаемых представлений и синтетических данных в отдельных задачах. Их совокупность обосновывает проверку сопряжения компонентов, но не устанавливает преимущество предобучения на динамике Drosophila для человеческого ECAP.
\item Оценка такого преимущества должна отделять необходимость временного модуля от эффекта источника предобучения. Для этого нужны простая модель кривой роста, человеческие модели сравнения и контроли структуры, динамики и адаптации при сопоставимой доступной информации.
\item Основная независимая единица целевой проверки --- участник. Средняя по участникам ошибка амплитуды, заранее установленные допуски и оценка неопределённости позволяют различать практически значимое преимущество, отсутствие существенного различия, ухудшение и недостаточную определённость результата.
\end{enumerate}

Научная проблема заключается в установлении условий, при которых коннектомное предобучение временного модуля даёт дополнительную ценность для физически обоснованного моделирования кривой роста ECAP. Её решение требует подготовки совместимых исходных наблюдений, разработки генератора динамики, определения переносимых параметров, отдельной проверки человеческого оператора и сравнительной оценки на независимых данных. Такая последовательность согласует обзор с целью и шестью задачами исследования, сформулированными во введении. Состав данных и проверяемые компоненты рассматриваются в следующей главе.
